\section*{Chapter \ref{ch:ll:protocols-ii-binary-exchange-protocols} --- Protocols II: Binary Exchange Protocols}

\begin{solution}{exo:ll:protocols-ii-binary-exchange-protocols:1}
The first packet carries 1\,000--1\,002, so the next expected number is 1\,003: messages 1\,003--1\,006 are missing (a gap of
four). The \omterm{def:ll:protocols-i-fix:session}{heartbeat}'s number is the next expected one, so the packet at 1\,007 carried three messages (1\,007--1\,009) and
nothing more was lost after it.
\end{solution}

\begin{solution}{exo:ll:protocols-ii-binary-exchange-protocols:2}
$\mathtt{0x000F42A4} = 1\,000\,100$ ten-thousandths: 100.01.
\end{solution}

\begin{solution}{exo:ll:protocols-ii-binary-exchange-protocols:3}
$0.002^2 = 4 \times 10^{-6}$ with independent lines. With the shared switch, $10^{-4} + (1 - 10^{-4}) \times 4 \times 10^{-6}
\approx 1.04 \times 10^{-4}$: the switch alone accounts for 96\% of the losses.
\end{solution}

\begin{solution}{exo:ll:protocols-ii-binary-exchange-protocols:4}
The bad state's stationary probability is $0.0005/0.1005 \approx 0.50\%$, so the line loses $0.8 \times 0.50\% \approx 0.40\%$
of packets; a bad period lasts $1/0.1 = 10$ packets on average. Losses therefore come in runs of several packets at once: a
retransmission request covers them, but at the moment they happen, which is often a burst of traffic, and one request may not
be enough.
\end{solution}

\begin{solution}{exo:ll:protocols-ii-binary-exchange-protocols:5}
One \omterm{def:ll:protocols-ii-binary-exchange-protocols:multicast}{multicast} packet reaches every subscriber at once, with no per-subscriber connection, acknowledgement or retransmission to
slow the publisher, and no subscriber can hold the others back. A subscriber gives up reliable, ordered delivery: it must
detect gaps from sequence numbers, arbitrate redundant lines and recover from retransmission or snapshots itself.
\end{solution}

\begin{solution}{exo:ll:protocols-ii-binary-exchange-protocols:6}
The message becomes 37 bytes of fields plus an 8-byte header, 45 against 36: 25\% more bandwidth and cache for Add Orders.
What it saves is the assembly of a six-byte integer; the measurement shows that this saving is not measurable, while the header
buys versioning and a length that lets old readers skip fields they do not know.
\end{solution}

\begin{solution}{exo:ll:protocols-ii-binary-exchange-protocols:7}
Add \texttt{soup\_client} and \texttt{soup\_server} to the protocols the generator walks, generating for the variable-length
messages (\texttt{U}, \texttt{S}, \texttt{+}) a reader that returns the payload as a view; the frame is a big-endian u16 length
that counts the type byte. The test reads \texttt{fixture\_soup.bin} frame by frame and compares each with
\texttt{fixture\_soup.csv}.
\end{solution}

\begin{solution}{exo:ll:protocols-ii-binary-exchange-protocols:8}
It does not arbitrate: while it reads one line, the other's packets are ignored, so each switch loses whatever arrived first on
the other line, and a gap on line A is only filled if line B's copies are still to come. The handler must read both lines all
the time and keep, per sequence number, the first copy from either.
\end{solution}

\begin{solution}{pb:ll:protocols-ii-binary-exchange-protocols:1}
\begin{enumerate}
\item $50\,000 \times 23\,400 = 1.17 \times 10^9$ packets; one line loses about 117\,000 of them.
\item $10^{-8}$, about 11.7 losses a session.
\item Because the losses were not independent: both lines were out together for 25 milliseconds (17 messages), and line B lost a
packet while line A was down.
\item One retransmission round trip, \qty{0.5}{\milli\second}, during which the book cannot be trusted.
\item $10^{-6} + (1 - 10^{-6}) \times 10^{-8} \approx 1.01 \times 10^{-6}$: about 1\,182 losses a session, a hundred times more.
\item About 0.40\% per line, and $1.6 \times 10^{-5}$ arbitrated if the lines are independent.
\item Losses arrive in runs at the busiest moments; the request, the answer and the replay happen while the market moves, and a
run longer than the retransmission limit needs a snapshot.
\item A shared switch, a shared network card or server, a shared cross-connect or carrier, shared power; it cannot remove the
publisher's own failures, which are common to both lines.
\item When the gap exceeds what the server keeps or serves (1\,000 messages per request, a window of 100\,000), when the
server is unreachable, or when the handler starts late.
\item Half an interval plus the snapshot on average, $0.5 + 0.02 = \qty{0.52}{\second}$; a full interval at worst,
\qty{1.02}{\second}.
\item About 26\,000 messages on average (51\,000 at worst), applied in about \qty{0.21}{\milli\second} (\qty{0.41}{\milli\second}).
\item Those the snapshot already reflects: sequence numbers up to the one the snapshot's begin marker gives (CME's
LastMsgSeqNumProcessed).
\item \textbf{Named result.} With independent losses of $10^{-4}$ per line the arbitrated stream loses $10^{-8}$ per packet,
about 12 packets a session at 50\,000 a second; a common cause of $10^{-6}$ makes it $1.01 \times 10^{-6}$, about 1\,180 a
session. Recovery from a one-second snapshot cycle takes \qty{0.52}{\second} on average and \qty{1.02}{\second} at worst,
against \qty{0.5}{\milli\second} for a retransmission.
\item Mark the book stale, stop quoting on it (or quote only defensively) and cancel what depends on it until it is rebuilt.
\item It cuts the average wait to \qty{50}{\milli\second}, at the price of ten times the snapshot bandwidth for every
subscriber, most of whom never need it.
\item Waiting for the slower line adds that line's delay to every message; the point of two lines is to take the faster copy.
\item Replay recorded lines with losses injected, one line at a time and both together (as the simulator does), and check that
the rebuilt book equals the book of a lossless run.
\item The line, the sequence range, the time, whether it was recovered and how, and how long the book was stale.
\item Almost nowhere in decoding: a few nanoseconds per message; the latency is in the network, the kernel and the book.
\item Protection from independent failures of one path, not from failures common to both.
\end{enumerate}
\end{solution}

\begin{solution}{iq:ll:protocols-ii-binary-exchange-protocols:1}
One packet reaches every subscriber at once, and the publisher never waits for anyone. Each packet carries the sequence number
of its first message and a count; a packet that starts above the next expected number reveals a gap, and \omterm{def:ll:protocols-i-fix:session}{heartbeats} carry the
next number so that losses are seen even when nothing trades.
\iqlookfor{fairness and scale; sequence numbers and \omterm{def:ll:protocols-i-fix:session}{heartbeats}.}
\end{solution}

\begin{solution}{iq:ll:protocols-ii-binary-exchange-protocols:2}
Keep the next expected sequence number. For each packet from either line: if it ends before that number, drop it; if it starts
after it, record a gap (lost on both) and start recovery; otherwise deliver the new messages and advance. Read both lines
without waiting for either; buffer what arrives during recovery.
\iqlookfor{one counter, duplicates dropped, gaps recovered.}
\end{solution}

\begin{solution}{iq:ll:protocols-ii-binary-exchange-protocols:3}
An object holding only a pointer to the message, whose accessors read each field from its fixed offset when asked. Nothing is
copied, nothing allocated, unread fields cost nothing, and there is no indirect call: the measured cost is half that of a
decoder that copies into a struct and calls a callback.
\iqlookfor{views over bytes and lazy field access.}
\end{solution}

\begin{solution}{iq:ll:protocols-ii-binary-exchange-protocols:4}
Subscribe to the incremental lines and buffer; wait for a snapshot of each instrument; load it; drop buffered messages up to
the sequence number the snapshot reflects; apply the rest in order; then process live. Verify with the snapshot's checksum where
there is one.
\iqlookfor{buffer, snapshot, sequence number alignment.}
\end{solution}

\begin{solution}{iq:ll:protocols-ii-binary-exchange-protocols:5}
Hardly: a byte swap is one instruction, and the chapter's measurement shows no difference between big- and little-endian
flyweights. What matters is not copying, not allocating, not calling through pointers, and fixed offsets that need no search.
\iqlookfor{measure; the costs are elsewhere.}
\end{solution}

\begin{solution}{iq:ll:protocols-ii-binary-exchange-protocols:6}
A common cause: shared switches, cards, cables, carriers or power; the publisher itself (both lines losing the same range at the
source); the receiving host (both lines' interrupts on one overloaded core, a full socket buffer, a slow handler). Compare gap
times and ranges across lines and with other subscribers, and check the host's drop counters.
\iqlookfor{independence as the assumption to test, including the receiving host.}
\end{solution}
